%% toc: P15 159 159.7,159.10
\section{Regular conditional distributions defined almost everywhere}
\label{sec:annex-p15-w02b-p15-008-ae-kernels}

Throughout, \(\mu\) is a regular conditional distribution for \(X\) given
\(\mathcal G\) in the sense of \cref{def:p15-w02b-p15-008-rcd}, with
exceptional set \(N\in\mathcal F\), \(P(N)=0\), as in
\cref{def:p15-w02b-p15-008-exceptional-set}. By the typing of \(\mu\),
every value \(\mu(\omega,A)\) lies in \([0,1]\), including for
\(\omega\in N\).

\begin{annexitem}{annex:p15-w02b-p15-008-ae-integration}{thm:p15-w02b-p15-008-ae-integration}{Indicators, simple functions, monotone limits and the signed case off one exceptional set}
\begin{proof}
\paragraph{Simple functions.}
Let \(\varphi=\sum_{j=1}^m c_j\mathbf1_{A_j}\) with \(c_j\in[0,\infty)\)
and \(A_j\in\mathcal S\), the representation being fixed. Define
\[
            g_\varphi(\omega)=\sum_{j=1}^m c_j\,\mu(\omega,A_j)
                                   \qquad(\omega\in\Omega).
\]
Each \(\mu(\cdot,A_j)\) is a version of \(P(X\in A_j\mid\mathcal G)\), hence
\(\mathcal G\)-measurable, so \(g_\varphi\) is \(\mathcal G\)-measurable and
takes values in \([0,\infty)\).
For \(\omega\notin N\), \(\mu(\omega,\cdot)\) is a probability measure and
linearity of its integral gives \(\int\mu(\omega,dx)\varphi(x)=g_\varphi(\omega)\).
For \(B\in\mathcal G\), the defining identity of a version of
\(P(X\in A_j\mid\mathcal G)=E(\mathbf1_{\{X\in A_j\}}\mid\mathcal G)\) in
\cref{def:p15-general-conditional-expectation} gives
\(\int_B\mu(\cdot,A_j)\,dP=P(B\cap\{X\in A_j\})\). Summing,
\[
 \int_Bg_\varphi\,dP=\sum_{j=1}^mc_jP(B\cap\{X\in A_j\})
                    =E\bigl(\varphi(X)\mathbf1_B\bigr).
 \tag{1}
\]
For \(\varphi=\mathbf1_A\) this is the definition itself.

\paragraph{Nonnegative functions.}
Let \(f:S\to[0,\infty]\) be measurable. For \(n\geq1\) put
\[
           f_n(x)=\min\bigl\{n,\;2^{-n}\lfloor2^nf(x)\rfloor\bigr\},
\]
with \(f_n(x)=n\) when \(f(x)=\infty\).
Then \(f_n\) takes only the finitely many values \(k2^{-n}\),
\(0\leq k\leq n2^n\), and \(\{f_n=k2^{-n}\}\) is a finite Boolean
combination of the sets \(\{f\geq i2^{-n}\}\in\mathcal S\); so
\(f_n=\sum_kk2^{-n}\mathbf1_{\{f_n=k2^{-n}\}}\) is a nonnegative simple
function with a fixed representation.
The sequence increases: \(2\lfloor a\rfloor\leq\lfloor2a\rfloor\) for real
\(a\geq0\), so \(2^{-n}\lfloor2^nt\rfloor\leq2^{-n-1}\lfloor2^{n+1}t\rfloor\),
and \(n<n+1\). It converges to \(f\) pointwise: if \(f(x)<\infty\), then
\(0\leq f(x)-2^{-n}\lfloor2^nf(x)\rfloor<2^{-n}\) and \(n>f(x)\) for large
\(n\); if \(f(x)=\infty\), then \(f_n(x)=n\to\infty\).

Write \(g_n=g_{f_n}\) and \(g_f=\limsup_ng_n\), a
\(\mathcal G\)-measurable function with values in \([0,\infty]\).
For \(\omega\notin N\), the simple-function step gives
\[
                  g_n(\omega)=\int\mu(\omega,dx)f_n(x).
\]
These numbers increase in \(n\),
and the monotone convergence theorem for the probability measure
\(\mu(\omega,\cdot)\) gives
\[
       g_f(\omega)=\lim_ng_n(\omega)=\int\mu(\omega,dx)f(x)
                                       \qquad(\omega\notin N).
\]
Let \(B\in\mathcal G\). The functions \(g_n\mathbf1_{B\setminus N}\) are
\(\mathcal F\)-measurable, because \(N\in\mathcal F\), and they increase
to \(g_f\mathbf1_{B\setminus N}\) at every point of \(\Omega\).
Since \(P(N)=0\), monotone convergence on \((\Omega,\mathcal F,P)\) and (1)
give
\[
 \int_Bg_f\,dP=\int_{B\setminus N}g_f\,dP
   =\lim_n\int_{B\setminus N}g_n\,dP
   =\lim_n\int_Bg_n\,dP
   =\lim_nE\bigl(f_n(X)\mathbf1_B\bigr).
\]
The variables \(f_n(X)\mathbf1_B\) increase pointwise to
\(f(X)\mathbf1_B\), so a second application of monotone convergence gives
\[
          \int_Bg_f\,dP=E\bigl(f(X)\mathbf1_B\bigr)
                          \qquad(B\in\mathcal G).
          \tag{2}
\]
Thus \(g_f\) is a nonnegative extended \(\mathcal G\)-measurable function
satisfying the all-event identity that characterizes the nonnegative
conditional expectation, and by the uniqueness clause of
\cref{thm:p15-conditional-limits} it is a version of
\(C_{\mathcal G}f(X)\). The set \(\Omega\setminus N\) was fixed before
\(f\) was chosen. This proves part~(i).

\paragraph{Integrable real functions.}
Let \(f\) be real and measurable with \(E|f(X)|<\infty\). Put
\(f^+=\max(f,0)\) and \(f^-=\max(-f,0)\), measurable with
\(|f|=f^++f^-\) and \(f=f^+-f^-\). Apply part~(i) to obtain
\(g_+=g_{f^+}\) and \(g_-=g_{f^-}\). Taking \(B=\Omega\) in (2),
\[
          \int g_\pm\,dP=E f^\pm(X)\leq E|f(X)|<\infty .
\]
Let \(M_f=\{g_+=\infty\}\cup\{g_-=\infty\}\). It belongs to \(\mathcal G\)
because \(g_\pm\) are \(\mathcal G\)-measurable, and \(P(M_f)=0\) because
a nonnegative function with finite integral is finite outside a null set.

For \(\omega\notin N\cup M_f\), part~(i) gives
\(\int\mu(\omega,dx)f^\pm(x)=g_\pm(\omega)<\infty\). Additivity of the
integral for nonnegative functions yields
\[
 \int\mu(\omega,dx)|f(x)|=g_+(\omega)+g_-(\omega)<\infty,
\]
so \(f\) is \(\mu(\omega,\cdot)\)-integrable and
\(\int\mu(\omega,dx)f(x)=g_+(\omega)-g_-(\omega)\).

Define \(Z_f=(g_+-g_-)\mathbf1_{\Omega\setminus M_f}\). It is real-valued
and \(\mathcal G\)-measurable, and \(|Z_f|\leq g_++g_-\), so
\(E|Z_f|\leq E|f(X)|<\infty\). For \(B\in\mathcal G\), using
\(P(M_f)=0\) and (2) for \(f^+\) and \(f^-\),
\[
 \int_BZ_f\,dP=\int_Bg_+\,dP-\int_Bg_-\,dP,
\]
and the right side equals
\(E(f^+(X)\mathbf1_B)-E(f^-(X)\mathbf1_B)=E(f(X)\mathbf1_B)\).
All event integrals of \(Z_f\) are defined, so \(Z_f\) is a version of
\(E(f(X)\mid\mathcal G)\) in the sense of
\cref{def:p15-general-conditional-expectation}. By the preceding paragraph,
\(Z_f(\omega)=\int\mu(\omega,dx)f(x)\) for every
\(\omega\notin N\cup M_f\). Since \(P(N\cup M_f)=0\), and every version of
\(E(f(X)\mid\mathcal G)\) agrees with \(Z_f\) almost surely by
\cref{thm:p15-conditional-existence}, the almost-sure formula of part~(ii)
follows.
\end{proof}
\end{annexitem}

\begin{annexitem}{annex:p15-w02b-p15-008-ae-reduction}{thm:p15-w02b-p15-008-ae-reduction}{A point-mass default on a conditioning-measurable null set}
\begin{proof}
For \(A\in\mathcal S\),
\[
 K(\cdot,A)=\mu(\cdot,A)\mathbf1_{\Omega\setminus M_0}
           +\mathbf1_A(s_0)\mathbf1_{M_0}
\]
is a sum of products of \(\mathcal G\)-measurable functions, because
\(\mu(\cdot,A)\) is \(\mathcal G\)-measurable and \(M_0\in\mathcal G\).
For \(\omega\in M_0\), \(K(\omega,\cdot)\) is the point mass at \(s_0\), a
probability measure. For \(\omega\notin M_0\) we have \(\omega\notin N\)
because \(N\subseteq M_0\), so \(K(\omega,\cdot)=\mu(\omega,\cdot)\) is a
probability measure. Hence \(K\) is a probability kernel from
\((\Omega,\mathcal G)\) to \((S,\mathcal S)\) in the sense of
\cref{def:p15-probability-kernel}.

For each \(A\), \(K(\cdot,A)\) and \(\mu(\cdot,A)\) differ at most on
\(M_0\), which is null. Hence they have the same integral over every
\(B\in\mathcal G\); since \(\mu(\cdot,A)\) is a version of
\(P(X\in A\mid\mathcal G)\), so is \(K(\cdot,A)\). This is the identity
required by \cref{def:p15-regular-conditional-law}.

For measurable \(f\geq0\), \cref{thm:p15-kernel-expectation} for \(K\)
and part~(i) of \cref{thm:p15-w02b-p15-008-ae-integration} for \(\mu\)
each give a version of \(C_{\mathcal G}f(X)\); the uniqueness clause of
\cref{thm:p15-conditional-limits} makes them equal almost surely. For real
\(f\) with \(E|f(X)|<\infty\), both theorems give versions of
\(E(f(X)\mid\mathcal G)\), which agree almost surely by
\cref{thm:p15-conditional-existence}. At \(\omega\notin M_0\) the measures
\(K(\omega,\cdot)\) and \(\mu(\omega,\cdot)\) are equal, so the two kernel
integrals of \(f\) are equal wherever one of them is defined.

If \(\mathcal G\) contains every \(P\)-null set of \(\mathcal F\), then
\(N\in\mathcal G\), because \(N\in\mathcal F\) and \(P(N)=0\), and
\(M_0=N\) satisfies the hypothesis.
\end{proof}
\end{annexitem}

\begin{annexitem}{annex:p15-w02b-p15-008-ae-counterexample}{thm:p15-w02b-p15-008-ae-counterexample}{The countable--cocountable model}
\begin{proof}
\paragraph{The spaces.}
\(P\) is a probability measure on \(\mathcal F\): it is the restriction of
Lebesgue measure to \([0,1]\), transported to subsets of \([0,2]\) by
intersection, and \(P([0,2])=\lambda([0,1])=1\). The same formula defines
\(\nu\) on \(\mathcal S=\mathcal F\), so \(\nu=P\).

\(\mathcal G\) is a sigma-algebra. It contains \(\emptyset\). It is closed
under complements because its defining condition is symmetric in a set
and its complement. Let \(A_1,A_2,\ldots\in\mathcal G\). If every
\(A_k\) is countable, so is their union. If some \(A_k\) has countable
complement, the complement of the union is contained in
\([0,2]\setminus A_k\) and is countable.
Every countable set is a countable union of closed singletons and is Borel,
and so is its complement; hence \(\mathcal G\subseteq\mathcal F\).
A countable set \(C\) has \(P(C)=0\) by countable additivity and
\(\lambda(\{t\})=0\). A set with countable complement then has
probability \(1\). In particular every \(M\in\mathcal G\) with \(P(M)=0\)
is countable.

The interval \((1,2]\) is uncountable: a countable set has Lebesgue
measure \(0\), while \(\lambda((1,2])=1\). For the same reason \([0,1]\)
is uncountable. So neither \([0,1]\) nor \((1,2]\) belongs to
\(\mathcal G\).

\paragraph{Part (i).}
Every value of \(\mu\) lies in \([0,1]\). Fix \(A\in\mathcal S\) and let
\(c_A=\nu(A)\). The constant function \(c_A\) is a version of
\(P(X\in A\mid\mathcal G)\). It is \(\mathcal G\)-measurable. If
\(B\in\mathcal G\) is countable, then \(\int_Bc_A\,dP=0\) and
\(P(B\cap\{X\in A\})\leq P(B)=0\). If \(B\) has countable complement, then
\(A\setminus B\subseteq\Omega\setminus B\) is countable and
\(P(A\setminus B)=0\), so
\[
 P(B\cap\{X\in A\})=P(A)=c_A=c_AP(B)=\int_Bc_A\,dP .
\]
Now compare \(\mu(\cdot,A)\) with \(c_A\). If \(A\) is not a singleton,
then \(\mu(\omega,A)=\nu(A)\) for every \(\omega\), so
\(\mu(\cdot,A)=c_A\). If \(A=\{a\}\), then \(\nu(\{a\})=0\) and
\[
 \mu(\cdot,\{a\})=\tfrac12\mathbf1_{\{a\}\cap(1,2]},
\]
which is \(\mathcal G\)-measurable because \(\{a\}\cap(1,2]\) is
countable. In both cases \(\mu(\cdot,A)\) is \(\mathcal G\)-measurable and
differs from \(c_A\) on at most one point, a \(P\)-null set. The two
functions have the same integral over every \(B\in\mathcal G\), so
\(\mu(\cdot,A)\) is a version of \(P(X\in A\mid\mathcal G)\). Condition~(i)
of \cref{def:p15-w02b-p15-008-rcd} holds.

For \(\omega\in[0,1]\) and every \(A\), the first case of the definition
of \(\mu\) does not apply, so \(\mu(\omega,\cdot)=\nu\), a probability
measure. For \(\omega\in(1,2]\), the set \([0,2]\setminus\{\omega\}\) is not
a singleton and does not meet \(\{\omega\}\), and
\[
 \mu(\omega,\{\omega\})+\mu(\omega,[0,2]\setminus\{\omega\})
   =\tfrac12+\lambda([0,1])=\tfrac32
   \neq1=\nu([0,2])=\mu(\omega,[0,2]).
\]
So \(\mu(\omega,\cdot)\) is not additive and is not a measure. The set of
arguments where the measure condition fails is therefore exactly
\(N=(1,2]\), a Borel set with \(P(N)=\lambda(\emptyset)=0\). This verifies
condition~(ii) with exceptional set \(N\).

\paragraph{Part (ii).}
If \(M\in\mathcal G\), \(P(M)=0\) and \(N\subseteq M\), then \(M\) is
countable by the paragraph on the spaces, and so is \(N\); this contradicts the
uncountability of \((1,2]\). Suppose now that \(K\) is a probability kernel
from \((\Omega,\mathcal G)\) to \((S,\mathcal S)\), that
\(M\in\mathcal G\), \(P(M)=0\), and that \(K(\omega,\cdot)=\mu(\omega,\cdot)\)
for all \(\omega\notin M\). For \(\omega\in N\), \(K(\omega,\cdot)\) is a
probability measure and \(\mu(\omega,\cdot)\) is not, so \(\omega\in M\).
Thus \(N\subseteq M\), which was just shown impossible.

\paragraph{Part (iii).}
For each \(A\), \(K_0(\cdot,A)=c_A\) is constant, hence
\(\mathcal G\)-measurable, and it is a version of
\(P(X\in A\mid\mathcal G)\) by part~(i). Every \(K_0(\omega,\cdot)=\nu\) is
a probability measure. So \(K_0\) is a regular conditional law in the sense
of \cref{def:p15-regular-conditional-law}. The comparison in part~(i)
shows that \(K_0(\cdot,A)=\mu(\cdot,A)\) except at most at the single point
\(a\) when \(A=\{a\}\) with \(a\in(1,2]\).

\paragraph{Part (iv).}
Let \(f\) be Borel with \(E|f(X)|=\int|f|\,d\nu<\infty\), and
\(c_f=\int f\,d\nu\). For every \(\omega\),
\(\int K_0(\omega,dx)|f(x)|=\int|f|\,d\nu<\infty\) and
\(\int K_0(\omega,dx)f(x)=c_f\). The absolute kernel integral is finite at
every argument, so the version supplied by
\cref{thm:p15-kernel-expectation} for \(K_0\) is the constant \(c_f\) on
all of \(\Omega\). For \(\omega\in[0,1]\),
\(\mu(\omega,\cdot)=\nu\), so \(\int\mu(\omega,dx)f(x)=c_f\).

Let \(h=c_f\mathbf1_{[0,1]}\). If \(c_f=0\), then \(h\equiv0\) is
\(\mathcal G\)-measurable. If \(c_f\neq0\), then
\(h^{-1}(\{c_f\})=[0,1]\notin\mathcal G\), so \(h\) is not
\(\mathcal G\)-measurable.
\end{proof}
\end{annexitem}
